\documentclass[11pt]{article}
\usepackage[margin=0.82in]{geometry}

% Typography and mathematics
\usepackage[T1]{fontenc}
\usepackage{lmodern}
\usepackage{microtype}
\usepackage{amsmath,amssymb}

% Layout and visual structure
\usepackage[dvipsnames]{xcolor}
\usepackage{booktabs}
\usepackage{tabularx}
\usepackage{array}
\usepackage{enumitem}
\usepackage{multicol}
\usepackage[most]{tcolorbox}
\usepackage{fancyhdr}
\usepackage{titlesec}
\usepackage{hyperref}

% Colour palette
\definecolor{PlanNavy}{HTML}{17324D}
\definecolor{PlanBlue}{HTML}{2F6690}
\definecolor{PlanTeal}{HTML}{3A7D78}
\definecolor{PlanGold}{HTML}{D89B35}
\definecolor{PlanRed}{HTML}{B64B4B}
\definecolor{PlanInk}{HTML}{263238}
\definecolor{PlanPaper}{HTML}{F5F8FA}
\definecolor{PlanLine}{HTML}{D7E1E8}

\hypersetup{
  colorlinks=true,
  linkcolor=PlanBlue,
  urlcolor=PlanTeal,
  pdftitle={Final Exam Study Plan}
}

% Paragraphs and lists
\setlength{\parindent}{0pt}
\setlength{\parskip}{0.55\baselineskip}
\setlist[itemize]{leftmargin=1.5em,itemsep=0.15em,topsep=0.25em}
\setlist[enumerate]{leftmargin=1.7em,itemsep=0.2em,topsep=0.25em}
\renewcommand{\arraystretch}{1.22}

% Headings and page furniture
\titleformat{\section}
  {\Large\bfseries\color{PlanNavy}}
  {}{0pt}{}
  [\vspace{0.15em}\color{PlanGold}\titlerule]
\titleformat{\subsection}
  {\large\bfseries\color{PlanBlue}}
  {}{0pt}{}
\titleformat{\subsubsection}
  {\normalsize\bfseries\color{PlanTeal}}
  {}{0pt}{}
\titlespacing*{\section}{0pt}{1.7em}{0.8em}
\titlespacing*{\subsection}{0pt}{1.2em}{0.45em}
\titlespacing*{\subsubsection}{0pt}{0.9em}{0.25em}

\pagestyle{fancy}
\fancyhf{}
\lhead{\small\color{PlanNavy}\textbf{Final Exam Study Plan}}
\rhead{\small\color{PlanBlue}ANN \textbullet{} CNN \textbullet{} Object Detection}
\cfoot{\small\color{PlanInk}\thepage}
\setlength{\headheight}{14pt}

% Reusable boxes
\newtcolorbox{focusbox}[1][]{
  enhanced,
  breakable,
  colback=PlanPaper,
  colframe=PlanBlue,
  boxrule=0.8pt,
  arc=2mm,
  left=3mm,right=3mm,top=2.2mm,bottom=2.2mm,
  borderline west={3pt}{0pt}{PlanBlue},
  title=#1,
  fonttitle=\bfseries,
  coltitle=PlanNavy
}

\newtcolorbox{prioritybox}[1]{
  enhanced,
  breakable,
  colback=PlanGold!8,
  colframe=PlanGold!80!black,
  boxrule=0.8pt,
  arc=2mm,
  left=3mm,right=3mm,top=2.2mm,bottom=2.2mm,
  title=#1,
  fonttitle=\bfseries,
  coltitle=PlanNavy
}

\newtcolorbox{formulabox}{
  enhanced,
  colback=PlanTeal!6,
  colframe=PlanTeal,
  boxrule=0.7pt,
  arc=1.5mm,
  left=2mm,right=2mm,top=1mm,bottom=1mm
}

\newtcolorbox{practicebox}[1]{
  enhanced,
  breakable,
  colback=white,
  colframe=PlanBlue!70,
  colbacktitle=PlanBlue!10,
  coltitle=PlanNavy,
  fonttitle=\bfseries,
  boxrule=0.7pt,
  arc=2mm,
  left=3mm,right=3mm,top=2mm,bottom=2mm,
  title={#1}
}

\newcommand{\highpriority}{\textcolor{PlanRed}{\textbf{Highest priority}}}
\newcommand{\corepriority}{\textcolor{PlanGold!70!black}{\textbf{Core priority}}}
\newcommand{\quickreview}{\textcolor{PlanTeal}{\textbf{Short review}}}
\newcommand{\checkbox}{\(\square\)}

\begin{document}

% Title block
\begin{tcolorbox}[
  enhanced,
  colback=PlanNavy,
  colframe=PlanNavy,
  arc=3mm,
  boxrule=0pt,
  left=5mm,right=5mm,top=5mm,bottom=5mm
]
  {\color{white}\fontsize{24}{28}\selectfont\bfseries Final Exam Study Plan}\par
  \vspace{0.35em}
  {\color{white!82}\large A practice-first roadmap for ANN, CNN, and Object Detection}\par
  \vspace{0.7em}
  {\color{PlanGold}\bfseries See the pattern \(\longrightarrow\) choose the rule \(\longrightarrow\) execute quickly}
\end{tcolorbox}

\begin{prioritybox}{Preparation strategy}
The course material points to a clear emphasis: the ANN work progresses from the chain rule to full matrix backpropagation; the CNN work is dominated by numerical calculations; and the object-detection material repeatedly returns to IoU, NMS, anchors, YOLO, R-CNN variants, and U-Net. Therefore, the final phase should \textbf{not} divide time equally across every topic. It should be driven by problem-solving and exam patterns, not by relearning theory.
\end{prioritybox}

\subsection{Plan at a glance}

\begin{center}
\begin{tabularx}{\textwidth}{>{\bfseries}l X >{\centering\arraybackslash}p{2.4cm} >{\centering\arraybackslash}p{2.6cm}}
\toprule
Area & Main emphasis & Target practice & Priority \\
\midrule
ANN & Forward/backward variations, mixed activations, dimensions, and updates & 12--15 problems & \corepriority \\
CNN & Shapes, manual convolution, pooling, parameter counts, and efficient blocks & 10--12 problems & \corepriority \\
Object Detection & IoU, thresholding, class-wise NMS, TP/FP, anchors, and YOLO & 8--10 problems & \highpriority \\
Mock Finals & Mixed, timed execution across all three areas & 3 papers & \highpriority \\
\bottomrule
\end{tabularx}
\end{center}

\begin{focusbox}[Exact order]
\begin{center}
\large\bfseries\color{PlanNavy}
ANN Variations
\(\boldsymbol{\longrightarrow}\)
CNN Numerical Drills
\(\boldsymbol{\longrightarrow}\)
Object Detection Drills
\(\boldsymbol{\longrightarrow}\)
3 Mock Finals
\end{center}
\end{focusbox}

% -----------------------------------------------------------------------------
\section{Part I: Artificial Neural Networks}

\begin{focusbox}[Objective]
Master every common forward/backward variation. The goal is to recognize the output-layer shortcut, propagate the error with the correct activation derivative, and produce gradients with the correct dimensions without re-deriving backpropagation from scratch.
\end{focusbox}

\subsection{Stage A: One hidden layer}

Practice all four base architectures:
\[
2\!\to\!2\!\to\!1,\qquad
3\!\to\!2\!\to\!1,\qquad
2\!\to\!3\!\to\!1,\qquad
3\!\to\!3\!\to\!1.
\]

Rotate through the following activation--loss combinations.

\begin{center}
\small
\begin{tabularx}{\textwidth}{>{\bfseries}p{0.82cm} p{3.3cm} p{3.08cm} X}
\toprule
Case & Hidden/output setup & Loss & Derivative to recognize immediately \\
\midrule
1 & ReLU \(\to\) Sigmoid & Binary cross-entropy & \(\delta_o=\hat y-y\) \\
2 & Sigmoid \(\to\) Sigmoid & Binary cross-entropy & \(\delta_o=\hat y-y\), with \(\sigma'(z)=a(1-a)\) in the hidden layer \\
3 & Tanh \(\to\) Sigmoid & Binary cross-entropy & \(\delta_o=\hat y-y\), with \(\tanh'(z)=1-a^2\) \\
4 & Sigmoid \(\to\) Sigmoid & Mean squared error & \(\delta_o=(\hat y-y)\hat y(1-\hat y)\) \\
5 & Tanh or ReLU \(\to\) Linear & Mean squared error & Regression shortcut: \(\delta_o=\hat y-y\) \\
6 & Any hidden layer \(\to\) 2-output Softmax & Categorical cross-entropy & \(\delta_o=\hat{\boldsymbol y}-\boldsymbol y\) \\
\bottomrule
\end{tabularx}
\end{center}

\begin{prioritybox}{Do not miss this variation}
Give extra attention to \textbf{Tanh hidden activation + Softmax output + categorical cross-entropy}. It combines a simple output error, \(\delta_o=\hat{\boldsymbol y}-\boldsymbol y\), with the hidden derivative \(1-a^2\), and it matches the style of the recent in-course question.
\end{prioritybox}

\subsection{Stage B: Multiple outputs}

Use the architectures
\[
3\!\to\!2\!\to\!2
\qquad\text{and}\qquad
2\!\to\!3\!\to\!2
\]
to separate two superficially similar cases.

\begin{center}
\begin{tabularx}{\textwidth}{>{\bfseries}p{3.1cm} X X}
\toprule
 & Softmax outputs & Independent sigmoid outputs \\
\midrule
Use case & Mutually exclusive classes & Multilabel classification \\
Output & \(\hat{\boldsymbol y}=\operatorname{softmax}(\boldsymbol z_o)\) & \(\hat y_i=\sigma(z_i)\) independently \\
Loss & \(L=-\sum_i y_i\ln\hat y_i\) & One binary cross-entropy term per output \\
Output error & \(\delta_o=\hat{\boldsymbol y}-\boldsymbol y\) & \(\delta_o=\hat{\boldsymbol y}-\boldsymbol y\) \\
Key distinction & Probabilities sum to \(1\) & Probabilities need not sum to \(1\) \\
\bottomrule
\end{tabularx}
\end{center}

\subsection{Stage C: Deep networks}

Progress through
\[
1\!\to\!1\!\to\!1,\qquad
2\!\to\!2\!\to\!1\!\to\!1,\qquad
2\!\to\!2\!\to\!2\!\to\!1,\qquad
2\!\to\!2\!\to\!2\!\to\!2.
\]

The recurring backward pattern is:

\begin{formulabox}
\[
\boxed{
\boldsymbol\delta_{\mathrm{current}}
=
\left(W_{\mathrm{next}}^{T}\boldsymbol\delta_{\mathrm{next}}\right)
\odot f'(\boldsymbol z_{\mathrm{current}})
}
\qquad
\boxed{
dW=\boldsymbol\delta\,\boldsymbol a_{\mathrm{previous}}^{T}
}
\]
\end{formulabox}

Deliberately include the following variations:

\begin{multicols}{2}
\begin{itemize}
  \item dead-ReLU cases;
  \item no-dead-ReLU cases;
  \item mixed ReLU + Sigmoid;
  \item mixed Tanh + ReLU;
  \item Softmax output;
  \item Sigmoid/BCE output;
  \item Linear or Sigmoid/MSE output.
\end{itemize}
\end{multicols}

\subsection{Stage D: Dimension and transpose drills}

Practice both conventions:
\[
\boldsymbol z=W\boldsymbol x+\boldsymbol b
\qquad\text{and}\qquad
\boldsymbol z=W^{T}\boldsymbol x+\boldsymbol b.
\]

Before calculating any values, write the shape of every object:
\[
\boldsymbol x:\ ?\!\times\!?,\qquad
W:\ ?\!\times\!?,\qquad
\boldsymbol z:\ ?\!\times\!?,\qquad
\boldsymbol\delta:\ ?\!\times\!?.
\]

\begin{prioritybox}{Non-negotiable dimension check}
\[
\boxed{\operatorname{shape}(dW)=\operatorname{shape}(W)}
\]
This single check catches most transpose mistakes, especially when the question uses the less familiar convention \(\boldsymbol z=W^T\boldsymbol x+\boldsymbol b\).
\end{prioritybox}

\subsection{Stage E: Six final ANN challenges}

\begin{enumerate}[label=\textbf{Challenge \Alph*:},leftmargin=*,align=left]
  \item \(3\!\to\!3\!\to\!2\), with Tanh \(\to\) Softmax \(\to\) categorical cross-entropy.
  \item \(3\!\to\!2\!\to\!1\), with Sigmoid \(\to\) Sigmoid \(\to\) MSE.
  \item \(2\!\to\!3\!\to\!2\!\to\!1\), with ReLU \(\to\) Tanh \(\to\) Sigmoid/BCE.
  \item A multi-neuron hidden layer containing exactly one dead ReLU.
  \item Weight matrices supplied using the transposed convention.
  \item One full gradient-descent update of \emph{every} weight and bias.
\end{enumerate}

% -----------------------------------------------------------------------------
\section{Part II: Convolutional Neural Networks}

\begin{focusbox}[Objective]
Treat CNN preparation as a numerical unit. Build speed in output dimensions, manual convolution, pooling, parameter counts, MACs, and efficient architecture blocks.
\end{focusbox}

\subsection{Type 1: Output shape}

Vary kernel size, stride, padding, and filter count on inputs such as \(64\times64\times128\). Include rectangular inputs and rectangular kernels.

\begin{formulabox}
\[
\boxed{
H_o=\left\lfloor\frac{H+2P_h-K_h}{S_h}\right\rfloor+1
}
\qquad
\boxed{
W_o=\left\lfloor\frac{W+2P_w-K_w}{S_w}\right\rfloor+1
}
\]
The output depth equals the number of filters.
\end{formulabox}

\subsection{Type 2: Manual convolution}

Use small \(4\times4\), \(5\times5\), and \(6\times6\) matrices with:

\begin{itemize}
  \item \(2\times2\) and \(3\times3\) filters;
  \item stride \(1\) and stride \(2\);
  \item padding \(0\) and padding \(1\);
  \item the full sequence \(\text{Convolution}\to\text{ReLU}\to\text{Pooling}\).
\end{itemize}

\subsection{Type 3: Pooling}

Practice both max and average pooling with \(2\times2,\ S=2\), \(2\times2,\ S=1\), and unusual dimensions where the floor operation changes the output size.

\subsection{Type 4: Full-network parameter counting}

For networks of the form
\[
\text{Input}\to\text{Conv}\to\text{Pool}\to\text{Conv}\to\text{Pool}
\to\text{Flatten}\to\text{FC}\to\text{FC},
\]
always calculate in this order:

\begin{enumerate}
  \item output shape after every convolution and pooling layer;
  \item parameters in each convolutional layer;
  \item flatten size;
  \item parameters in each fully connected layer;
  \item total trainable parameters.
\end{enumerate}

\begin{formulabox}
\[
\boxed{P_{\mathrm{conv}}=(K_hK_wC_{\mathrm{in}}+1)C_{\mathrm{out}}}
\qquad
\boxed{P_{\mathrm{FC}}=(n_{\mathrm{in}}+1)n_{\mathrm{out}}}
\]
\end{formulabox}

Complete at least \textbf{six} full-network parameter-count questions.

\subsection{Type 5: Standard vs. depthwise-separable convolution}

For an output feature map of spatial size \(H\times W\), input depth \(M\), output depth \(N\), and a \(K\times K\) kernel:

\begin{center}
\begin{tabularx}{\textwidth}{>{\bfseries}p{3.2cm} X X}
\toprule
Operation & Parameters & Multiply--accumulates (MACs) \\
\midrule
Standard convolution & \(K^2MN\) & \(HWK^2MN\) \\
Depthwise convolution & \(K^2M\) & \(HWK^2M\) \\
Pointwise convolution & \(MN\) & \(HWMN\) \\
Depthwise separable & \(K^2M+MN\) & \(HW(K^2M+MN)\) \\
\bottomrule
\end{tabularx}
\end{center}

Calculate the total parameters, total MACs, reduction factor, and percentage saving for both \(3\times3\) and \(5\times5\) kernels and several choices of \((M,N,H,W)\). These formulas count weights only; add biases only when the question explicitly requires them.

\subsection{Architecture-based numericals}

\begin{center}
\small
\begin{tabularx}{\textwidth}{>{\bfseries}p{3.2cm} X}
\toprule
Architecture & Numerical focus \\
\midrule
MobileNetV1 & Depthwise + pointwise parameters and MACs. \\
MobileNetV2 & Expanded channels, expand--depthwise--project parameters, output size, and whether a residual connection is allowed. \\
Inception & Branch output shapes, concatenated depth, total parameters, and savings from a \(1\times1\) bottleneck. \\
SqueezeNet & Fire-module parameters and output channels for \(1\times1\) squeeze and \(1\times1/3\times3\) expand branches. \\
ShuffleNet / ResNeXt & Grouped-convolution parameters: \(P=K^2C_{\mathrm{in}}C_{\mathrm{out}}/g\). \\
SENet & Squeeze-and-excitation parameters for \(C\to C/r\to C\). \\
DenseNet & Channel growth: \(C_{\mathrm{final}}=C_0+Lk\). \\
ConvNeXt & One block consisting of \(7\times7\) depthwise convolution followed by \(C\to4C\to C\). \\
\bottomrule
\end{tabularx}
\end{center}

\begin{prioritybox}{CNN emphasis}
Spend more time calculating than memorizing architectural history. Every architecture above should become a short shape, parameter, or MAC problem.
\end{prioritybox}

% -----------------------------------------------------------------------------
\section{Part III: Object Detection}

\begin{focusbox}[Objective --- \highpriority]
Make the full numerical pipeline automatic: confidence filtering \(\to\) class-wise NMS \(\to\) surviving predictions \(\to\) TP/FP classification. Add anchor assignment and YOLO output dimensions once that pipeline is secure.
\end{focusbox}

\subsection{Type 1: Intersection over Union}

IoU measures the overlap between two bounding boxes relative to their combined area. A threshold of \(0.5\) is commonly used as a criterion for sufficiently overlapping localization.

\begin{formulabox}
\[
\boxed{
\operatorname{IoU}(A,B)
=\frac{\text{intersection area}}{\text{union area}}
}
\qquad
\boxed{A_U=A_A+A_B-A_I}
\]
\end{formulabox}

\subsubsection{The six-step recipe}

For
\[
A=(x_{1A},y_{1A},x_{2A},y_{2A}),
\qquad
B=(x_{1B},y_{1B},x_{2B},y_{2B}),
\]
use the same sequence every time.

\begin{enumerate}[label=\textbf{Step \arabic*:},leftmargin=*,align=left]
  \item Find the intersection coordinates:
  \[
  x_L=\max(x_{1A},x_{1B}),\qquad
  y_T=\max(y_{1A},y_{1B}),
  \]
  \[
  x_R=\min(x_{2A},x_{2B}),\qquad
  y_B=\min(y_{2A},y_{2B}).
  \]

  \item Find the valid intersection width and height:
  \[
  \boxed{W_I=\max(0,x_R-x_L)},
  \qquad
  \boxed{H_I=\max(0,y_B-y_T)}.
  \]

  \item Calculate the intersection area:
  \[
  A_I=W_IH_I.
  \]

  \item Calculate the individual box areas:
  \[
  A_A=(x_{2A}-x_{1A})(y_{2A}-y_{1A}),
  \]
  \[
  A_B=(x_{2B}-x_{1B})(y_{2B}-y_{1B}).
  \]

  \item Calculate the union:
  \[
  \boxed{A_U=A_A+A_B-A_I}.
  \]

  \item Divide and compare with the required threshold:
  \[
  \boxed{\operatorname{IoU}(A,B)=\frac{A_I}{A_U}}.
  \]
\end{enumerate}

\begin{prioritybox}{Most important coordinate rule}
\[
\boxed{
\text{Intersection uses MAX for starting coordinates and MIN for ending coordinates.}
}
\]
The outer \(\max(0,\cdot)\) in \(W_I\) and \(H_I\) prevents a negative width or height when the boxes do not overlap.
\end{prioritybox}

\subsubsection{Worked example: partial overlap}

\begin{practicebox}{Example 1 --- \(A=(20,20,80,80)\), \(B=(40,30,100,90)\)}
The intersection coordinates are
\[
\begin{aligned}
x_L&=\max(20,40)=40, & y_T&=\max(20,30)=30,\\
x_R&=\min(80,100)=80, & y_B&=\min(80,90)=80.
\end{aligned}
\]
Thus, the overlapping rectangle runs from \((40,30)\) to \((80,80)\), giving
\[
W_I=80-40=40,
\qquad
H_I=80-30=50,
\qquad
\boxed{A_I=40\times50=2000}.
\]
The individual areas are
\[
A_A=(80-20)(80-20)=60\times60=3600,
\]
\[
A_B=(100-40)(90-30)=60\times60=3600.
\]
Therefore,
\[
A_U=3600+3600-2000=5200,
\]
and
\[
\boxed{
\operatorname{IoU}(A,B)=\frac{2000}{5200}\approx0.3846\approx0.385
}.
\]
Since \(0.385<0.5\), the boxes are \textbf{not} sufficiently overlapping at a threshold of \(0.5\). In same-class NMS with \(T=0.5\), the lower-confidence box would not be suppressed solely because of this IoU.
\end{practicebox}

\begin{practicebox}{Problem 1 --- Partial overlap}
\[
A=(10,10,50,50),
\qquad
B=(30,20,70,60).
\]
The intersection is bounded by
\[
x_L=30,\quad y_T=20,\quad x_R=50,\quad y_B=50,
\]
so
\[
W_I=20,\qquad H_I=30,
\qquad \boxed{A_I=20\times30=600}.
\]
Both boxes have area \(40\times40=1600\). Hence,
\[
A_U=1600+1600-600=2600,
\qquad
\boxed{\operatorname{IoU}=\frac{600}{2600}\approx0.231}.
\]
At a threshold of \(0.5\), there is \textbf{not enough overlap}.
\end{practicebox}

\begin{practicebox}{Problem 2 --- One box completely inside another}
\[
A=(20,20,100,100),
\qquad
B=(40,40,80,80).
\]
Because \(B\) lies completely inside \(A\), its full area is the intersection:
\[
A_A=(100-20)^2=80^2=6400,
\qquad
A_B=(80-40)^2=40^2=1600,
\]
\[
A_I=1600,
\qquad
A_U=6400+1600-1600=6400.
\]
Therefore,
\[
\boxed{\operatorname{IoU}=\frac{1600}{6400}=0.25}.
\]
\textbf{Key insight:} even though \(B\) is entirely inside \(A\), the IoU is only \(0.25\) because the union is the whole larger box.
\end{practicebox}

\begin{practicebox}{Problem 3 --- Strong overlap}
\[
A=(20,20,80,80),
\qquad
B=(25,25,85,85).
\]
The intersection runs from \((25,25)\) to \((80,80)\), so
\[
W_I=H_I=55,
\qquad
\boxed{A_I=55^2=3025}.
\]
Each box has area \(60\times60=3600\). Therefore,
\[
A_U=3600+3600-3025=4175,
\]
\[
\boxed{\operatorname{IoU}=\frac{3025}{4175}\approx0.725}.
\]
Since \(0.725\geq0.5\), this is \textbf{strong overlap}. In same-class NMS with \(T=0.5\), the lower-confidence prediction would normally be suppressed.
\end{practicebox}

\begin{practicebox}{Problem 4 --- No overlap}
\[
A=(10,10,40,40),
\qquad
B=(60,60,100,100).
\]
Here,
\[
x_L=\max(10,60)=60,
\qquad
x_R=\min(40,100)=40.
\]
Because \(x_R<x_L\), there is no valid intersection width:
\[
W_I=\max(0,40-60)=0,
\qquad
A_I=0.
\]
Thus,
\[
\boxed{\operatorname{IoU}=0}.
\]
\textbf{Exam shortcut:} if the boxes clearly do not overlap, write \(\operatorname{IoU}=0\) immediately.
\end{practicebox}

\begin{practicebox}{Problem 5 --- Boxes touching at an edge}
\[
A=(10,10,50,50),
\qquad
B=(50,10,90,50).
\]
The boxes meet at \(x=50\), but
\[
x_L=\max(10,50)=50,
\qquad
x_R=\min(50,90)=50,
\]
so
\[
W_I=50-50=0.
\]
Although the vertical overlap is \(50-10=40\), the intersection area is
\[
A_I=0\times40=0,
\qquad
\boxed{\operatorname{IoU}=0}.
\]
\textbf{Important trap:} boxes that touch only along an edge or at one corner have zero intersection area and therefore zero IoU.
\end{practicebox}

\begin{practicebox}{Bonus Problem 6 --- Very different box sizes}
\[
A=(0,0,100,100),
\qquad
B=(25,25,75,75).
\]
Box \(B\) is completely inside \(A\). The areas are
\[
A_A=100\times100=10000,
\qquad
A_B=50\times50=2500.
\]
Consequently,
\[
A_I=2500,
\qquad
A_U=10000,
\qquad
\boxed{\operatorname{IoU}=\frac{2500}{10000}=0.25}.
\]
Containment does not guarantee a high IoU: the relative box sizes still matter.
\end{practicebox}

\subsubsection{Final IoU pattern summary}

\begin{center}
\begin{tabularx}{\textwidth}{>{\bfseries}X X}
\toprule
Situation & Expected IoU behavior \\
\midrule
Identical boxes & \(1\) \\
Strong overlap & Close to \(1\) \\
Partial overlap & Strictly between \(0\) and \(1\) \\
Small box completely inside a large box & Can still be low \\
No overlap & \(0\) \\
Touching only at an edge or corner & \(0\) \\
\bottomrule
\end{tabularx}
\end{center}

\begin{focusbox}[Fast exam pattern]
\[
\boxed{
x_L=\max(x_{1A},x_{1B}),\quad
x_R=\min(x_{2A},x_{2B})
}
\]
\[
\boxed{
y_T=\max(y_{1A},y_{1B}),\quad
y_B=\min(y_{2A},y_{2B})
}
\]
\[
\boxed{
W_I=\max(0,x_R-x_L),\quad
H_I=\max(0,y_B-y_T)
}
\]
Once this coordinate pattern is automatic, the remaining work is only area arithmetic.
\end{focusbox}

\subsection{Type 2: Non-maximum suppression}

Once IoU is automatic, NMS becomes a repeated selection-and-removal procedure. Unless stated otherwise, all boxes in the following examples predict the \textbf{same class}.

\begin{formulabox}
\[
\boxed{\text{The highest-confidence box survives.}}
\]
\[
\boxed{
\text{Remove a lower-confidence box when }
\operatorname{IoU}\geq T_{\mathrm{NMS}}.
}
\]
\end{formulabox}

\subsubsection{Core NMS procedure}

\begin{enumerate}[label=\textbf{Step \arabic*:},leftmargin=*,align=left]
  \item Remove predictions below the confidence threshold, if one is given.
  \item Sort the remaining predictions from highest to lowest confidence.
  \item Select and keep the highest-confidence prediction.
  \item Compare its IoU with every lower-confidence prediction of the same class.
  \item Suppress predictions whose IoU is at least \(T_{\mathrm{NMS}}\).
  \item Select the next-highest surviving prediction and repeat.
\end{enumerate}

\begin{prioritybox}{The iterative rule}
Do not compare every prediction only with the original highest-confidence box. After the first round, continue from the \textbf{next-highest box that survived}.
\end{prioritybox}

\begin{practicebox}{Problem 1 --- Basic NMS}
The same-class confidence scores are
\[
A:0.92,\qquad B:0.85,\qquad C:0.70,\qquad D:0.60,
\]
with
\[
\operatorname{IoU}(A,B)=0.72,\quad
\operatorname{IoU}(A,C)=0.35,\quad
\operatorname{IoU}(A,D)=0.10,
\]
\[
\operatorname{IoU}(C,D)=0.62,
\qquad
T_{\mathrm{NMS}}=0.5.
\]

\textbf{Round 1:} \(A\) has the highest confidence, so keep \(A\). Then compare the remaining boxes with \(A\):
\[
\begin{array}{ccl}
B: & 0.72\geq0.5 & \Rightarrow \text{suppress }B,\\
C: & 0.35<0.5 & \Rightarrow \text{keep }C\text{ for now},\\
D: & 0.10<0.5 & \Rightarrow \text{keep }D\text{ for now}.
\end{array}
\]
The remaining candidates are \(C\) and \(D\).

\textbf{Round 2:} \(C\) has the larger remaining confidence, so keep \(C\). Since
\[
\operatorname{IoU}(C,D)=0.62\geq0.5,
\]
\(D\) is suppressed. Therefore,
\[
\boxed{\text{Final survivors}=\{A,C\}}.
\]
\end{practicebox}

\begin{practicebox}{Problem 2 --- Change the NMS threshold}
Use exactly the same predictions and IoUs, but set
\[
T_{\mathrm{NMS}}=0.7.
\]
Keep \(A\) first. Then
\[
\operatorname{IoU}(A,B)=0.72\geq0.7
\quad\Rightarrow\quad B\text{ is suppressed},
\]
while
\[
\operatorname{IoU}(A,C)=0.35<0.7,
\qquad
\operatorname{IoU}(A,D)=0.10<0.7,
\]
so \(C\) and \(D\) remain. Keep \(C\) next. Now
\[
\operatorname{IoU}(C,D)=0.62<0.7,
\]
so \(D\) is not suppressed. Hence,
\[
\boxed{\text{Final survivors}=\{A,C,D\}}.
\]

\begin{center}
\begin{tabularx}{0.75\textwidth}{>{\bfseries}X >{\centering\arraybackslash}X}
\toprule
NMS threshold & Final survivors \\
\midrule
\(T_{\mathrm{NMS}}=0.5\) & \(\{A,C\}\) \\
\(T_{\mathrm{NMS}}=0.7\) & \(\{A,C,D\}\) \\
\bottomrule
\end{tabularx}
\end{center}

\[
\boxed{
T_{\mathrm{NMS}}\uparrow
\quad\Longrightarrow\quad
\text{less aggressive suppression}
}
\]
A higher threshold requires greater overlap before removal, so more overlapping predictions tend to survive.
\end{practicebox}

\subsection{Type 3: Confidence threshold + NMS}

The order of operations is essential:
\[
\boxed{
\text{confidence filtering first}
\quad\Longrightarrow\quad
\text{NMS second}
}
\]

\begin{practicebox}{Problem 3 --- Confidence filtering followed by NMS}
The predictions are

\begin{center}
\begin{tabular}{cc}
\toprule
\textbf{Box} & \textbf{Confidence} \\
\midrule
\(A\) & \(0.91\) \\
\(B\) & \(0.83\) \\
\(C\) & \(0.68\) \\
\(D\) & \(0.57\) \\
\(E\) & \(0.42\) \\
\bottomrule
\end{tabular}
\end{center}

Use
\[
T_{\mathrm{conf}}=0.6,
\qquad
T_{\mathrm{NMS}}=0.5,
\]
with
\[
\operatorname{IoU}(A,B)=0.66,
\qquad
\operatorname{IoU}(A,C)=0.20,
\qquad
\operatorname{IoU}(B,C)=0.25.
\]

\textbf{Step 1 --- Confidence filtering.} Remove every score below \(0.6\):
\[
D:0.57\to\text{remove},
\qquad
E:0.42\to\text{remove}.
\]
The candidates entering NMS are
\[
\boxed{\{A,B,C\}}.
\]

\textbf{Step 2 --- NMS.} Keep \(A\), the highest-confidence prediction. Then
\[
\operatorname{IoU}(A,B)=0.66\geq0.5
\quad\Rightarrow\quad B\text{ is suppressed},
\]
but
\[
\operatorname{IoU}(A,C)=0.20<0.5
\quad\Rightarrow\quad C\text{ survives}.
\]
Keep \(C\). Therefore,
\[
\boxed{\text{Final survivors}=\{A,C\}}.
\]
This matches the in-course pattern exactly: low-confidence removal occurs before any IoU-based suppression.
\end{practicebox}

\begin{practicebox}{Problem 4 --- The second kept box suppresses another}
Suppose
\[
A:0.95,\qquad B:0.88,\qquad C:0.81,\qquad D:0.74,
\]
and
\[
\operatorname{IoU}(A,B)=0.65,\quad
\operatorname{IoU}(A,C)=0.30,\quad
\operatorname{IoU}(A,D)=0.20,
\]
\[
\operatorname{IoU}(C,D)=0.75,
\qquad
T_{\mathrm{NMS}}=0.5.
\]

\textbf{First round:} keep \(A\). Suppress \(B\) because \(0.65\geq0.5\), while \(C\) and \(D\) survive because their IoUs with \(A\) are below the threshold.

\textbf{Second round:} \(C\) is now the highest-confidence survivor, so keep \(C\). Since
\[
\operatorname{IoU}(C,D)=0.75\geq0.5,
\]
\(D\) is suppressed. Thus,
\[
\boxed{\text{Final survivors}=\{A,C\}}.
\]
\textbf{Lesson:} NMS must continue from the next-highest surviving prediction; it does not stop after comparisons with \(A\).
\end{practicebox}

\begin{practicebox}{Problem 5 --- Actual coordinates + NMS}
Consider three same-class predictions:
\[
A=(10,10,50,50),\quad s_A=0.95,
\]
\[
B=(15,15,55,55),\quad s_B=0.85,
\]
\[
C=(70,70,110,110),\quad s_C=0.75,
\qquad
T_{\mathrm{NMS}}=0.5.
\]

Boxes \(A\) and \(B\) each have area
\[
40\times40=1600.
\]
Their intersection runs from \((15,15)\) to \((50,50)\), so
\[
W_I=H_I=35,
\qquad
A_I=35^2=1225.
\]
The union is
\[
A_U=1600+1600-1225=1975,
\]
giving
\[
\boxed{
\operatorname{IoU}(A,B)=\frac{1225}{1975}\approx0.620
}.
\]
Keep \(A\), then suppress \(B\) because \(0.620\geq0.5\). Box \(C\) is spatially separate from \(A\), so
\[
\operatorname{IoU}(A,C)=0
\]
and \(C\) survives. Therefore,
\[
\boxed{\text{Final survivors}=\{A,C\}}.
\]
\end{practicebox}

\subsubsection{NMS exam recipe}

\begin{focusbox}[Write this sequence before calculating]
\begin{enumerate}[label=\textbf{\arabic*.},leftmargin=*,align=left,nosep]
  \item \textbf{Confidence filter:} \(s<T_{\mathrm{conf}}\Rightarrow\text{remove}\).
  \item \textbf{Sort:} arrange the remaining boxes from highest to lowest confidence.
  \item \textbf{Pick:} keep the highest-confidence box.
  \item \textbf{Compare:} if \(\operatorname{IoU}\geq T_{\mathrm{NMS}}\), suppress the lower-confidence box.
  \item \textbf{Repeat:} continue from the next-highest surviving box.
\end{enumerate}
\end{focusbox}

\begin{prioritybox}{The two threshold effects}
\[
\boxed{
T_{\mathrm{conf}}\uparrow
\quad\Longrightarrow\quad
\text{fewer boxes enter NMS}
}
\]
\[
\boxed{
T_{\mathrm{NMS}}\uparrow
\quad\Longrightarrow\quad
\text{more overlapping boxes survive NMS}
}
\]
The confidence threshold controls the \emph{initial candidate set}; the NMS threshold controls how much overlap is tolerated among those candidates.
\end{prioritybox}

\subsection{Type 4: TP, FP, and FN after NMS}

After NMS produces the final prediction set, evaluation compares each prediction with the ground-truth objects of the \textbf{same class}. A common matching threshold is \(T_{\mathrm{eval}}=0.5\).

\begin{formulabox}
\[
\boxed{
\operatorname{IoU}(\text{prediction},\text{ground truth})\geq0.5
\quad\Longrightarrow\quad
\mathrm{TP}
}
\]
\[
\boxed{
\operatorname{IoU}<0.5
\quad\Longrightarrow\quad
\mathrm{FP}
}
\]
\end{formulabox}

\begin{prioritybox}{One-to-one matching rule}
A ground-truth object can be matched only once. If two predictions both overlap the same ground truth sufficiently, process them in descending confidence order: the highest-confidence valid match becomes the true positive, while the duplicate becomes a false positive.
\[
\boxed{\text{One ground truth cannot produce two true positives.}}
\]
\end{prioritybox}

\begin{practicebox}{Problem 1 --- One ground truth}
Let
\[
G=(20,20,80,80),
\]
and suppose the following predictions remain after NMS:
\[
A=(25,25,75,75),\quad s_A=0.90,
\]
\[
B=(100,100,150,150),\quad s_B=0.70.
\]
Use \(T_{\mathrm{eval}}=0.5\).

The ground-truth area is
\[
A_G=60\times60=3600.
\]
Prediction \(A\) lies completely inside \(G\), and its area is
\[
A_A=50\times50=2500.
\]
Therefore,
\[
A_I=2500,
\qquad
A_U=3600+2500-2500=3600,
\]
so
\[
\boxed{
\operatorname{IoU}(A,G)=\frac{2500}{3600}\approx0.694
}.
\]
Since \(0.694\geq0.5\),
\[
\boxed{A=\mathrm{TP}}.
\]

Prediction \(B\) does not overlap \(G\), so
\[
\operatorname{IoU}(B,G)=0
\qquad\Longrightarrow\qquad
\boxed{B=\mathrm{FP}}.
\]
The final count is
\[
\boxed{\mathrm{TP}=1,\qquad \mathrm{FP}=1}.
\]
\end{practicebox}

\begin{practicebox}{Problem 2 --- Two predictions around one object}
Consider one ground truth and two predictions:
\[
G=(10,10,50,50),
\]
\[
A=(12,12,48,48),\quad s_A=0.92,
\qquad
B=(14,14,46,46),\quad s_B=0.80.
\]
Suppose both duplicates are still present. Since each lies inside \(G\),
\[
\operatorname{IoU}(A,G)=\frac{36^2}{40^2}
=\frac{1296}{1600}=0.81,
\]
\[
\operatorname{IoU}(B,G)=\frac{32^2}{40^2}
=\frac{1024}{1600}=0.64.
\]
Both exceed \(0.5\), but both refer to the same object. Process the predictions by confidence: \(A\) matches \(G\) first, leaving no unmatched ground truth for \(B\). Thus,
\[
\boxed{A=\mathrm{TP}},
\qquad
\boxed{B=\mathrm{FP}\text{ (duplicate)}}.
\]
This is why duplicate detections reduce precision even when their localization is individually good.
\end{practicebox}

\begin{practicebox}{Problem 3 --- Two ground truths}
Let
\[
G_1=(10,10,50,50),
\qquad
G_2=(80,80,120,120),
\]
with predictions
\[
A=(12,12,48,48),\quad s_A=0.95,
\]
\[
B=(78,78,122,122),\quad s_B=0.88,
\]
\[
C=(30,70,60,100),\quad s_C=0.65.
\]

Prediction \(A\) lies inside \(G_1\), giving
\[
\operatorname{IoU}(A,G_1)=\frac{36^2}{40^2}
=\frac{1296}{1600}=0.81,
\]
so \(A\) is a true positive matched to \(G_1\).

Prediction \(B\) has area \(44^2=1936\), while its intersection with \(G_2\) is \(40^2=1600\). Hence,
\[
\operatorname{IoU}(B,G_2)
=\frac{1600}{1936+1600-1600}
=\frac{1600}{1936}
\approx0.826,
\]
so \(B\) is a true positive matched to \(G_2\).

Prediction \(C\) does not overlap either ground truth sufficiently, so it is a false positive. Therefore,
\[
\boxed{A=\mathrm{TP},\qquad B=\mathrm{TP},\qquad C=\mathrm{FP}},
\]
\[
\boxed{\mathrm{TP}=2,\qquad \mathrm{FP}=1}.
\]
\end{practicebox}

\begin{practicebox}{Problem 4 --- A missed ground truth}
Let
\[
G_1=(10,10,50,50),
\qquad
G_2=(70,70,110,110),
\]
and
\[
A=(12,12,48,48),\quad s_A=0.93,
\qquad
B=(20,70,50,100),\quad s_B=0.72.
\]
Prediction \(A\) matches \(G_1\), so
\[
\boxed{A=\mathrm{TP}}.
\]
Prediction \(B\) does not match \(G_2\) sufficiently, so
\[
\boxed{B=\mathrm{FP}}.
\]
Ground truth \(G_2\) has no valid prediction left to match it. Therefore,
\[
\boxed{G_2=\mathrm{FN}},
\]
and the final count is
\[
\boxed{\mathrm{TP}=1,\qquad \mathrm{FP}=1,\qquad \mathrm{FN}=1}.
\]
This is the first place a \textbf{false negative} appears: it is a real object that the detector failed to find.
\end{practicebox}

\begin{practicebox}{Problem 5 --- Precision and recall}
Suppose evaluation gives
\[
\mathrm{TP}=8,
\qquad
\mathrm{FP}=2,
\qquad
\mathrm{FN}=4.
\]
Then
\[
\operatorname{Precision}
=\frac{\mathrm{TP}}{\mathrm{TP}+\mathrm{FP}}
=\frac{8}{8+2}
=\boxed{0.8},
\]
and
\[
\operatorname{Recall}
=\frac{\mathrm{TP}}{\mathrm{TP}+\mathrm{FN}}
=\frac{8}{8+4}
\approx\boxed{0.667}.
\]

\begin{center}
\begin{tabularx}{\textwidth}{>{\bfseries}p{2.7cm} X}
\toprule
Metric & Interpretation \\
\midrule
Precision & Among the objects predicted by the model, how many predictions were correct? \\
Recall & Among the real objects in the image or dataset, how many did the model find? \\
\bottomrule
\end{tabularx}
\end{center}
\end{practicebox}

\subsubsection{Evaluation exam pattern}

\begin{focusbox}[Apply predictions in confidence order]
\begin{enumerate}[label=\textbf{\arabic*.},leftmargin=*,align=left,nosep]
  \item Sort the predictions from highest to lowest confidence.
  \item For each prediction, find the best \emph{unmatched} ground truth of the same class.
  \item If the best IoU satisfies \(\operatorname{IoU}\geq T_{\mathrm{eval}}\), label the prediction \(\mathrm{TP}\) and mark that ground truth as matched.
  \item If no valid unmatched ground truth exists, label the prediction \(\mathrm{FP}\).
  \item After all predictions are processed, every unmatched ground truth is an \(\mathrm{FN}\).
\end{enumerate}
\end{focusbox}

\begin{practicebox}{Harder practice --- Duplicates, two objects, and an extra box}
Use
\[
G_1=(10,10,50,50),
\qquad
G_2=(60,60,100,100),
\]
and the confidence-ordered predictions
\[
A=(12,12,48,48),\quad s_A=0.95,
\]
\[
B=(14,14,46,46),\quad s_B=0.85,
\]
\[
C=(62,62,98,98),\quad s_C=0.80,
\]
\[
D=(110,110,150,150),\quad s_D=0.60.
\]
At \(T_{\mathrm{eval}}=0.5\):
\begin{itemize}[nosep]
  \item \(A\) matches \(G_1\), so \(A=\mathrm{TP}\).
  \item \(B\) also overlaps \(G_1\), but \(G_1\) is already matched; hence \(B=\mathrm{FP}\).
  \item \(C\) matches \(G_2\), so \(C=\mathrm{TP}\).
  \item \(D\) matches neither ground truth, so \(D=\mathrm{FP}\).
\end{itemize}
Both ground truths are matched, so there are no false negatives. The result is
\[
\boxed{\mathrm{TP}=2,\qquad \mathrm{FP}=2,\qquad \mathrm{FN}=0}.
\]
\end{practicebox}

\begin{prioritybox}{Evaluation summary}
\begin{center}
\begin{tabularx}{0.9\textwidth}{>{\bfseries}p{1.6cm} X}
\toprule
Label & Meaning \\
\midrule
TP & A prediction correctly claims an unmatched ground-truth object. \\
FP & A prediction is poorly localized, has the wrong class, is a duplicate, or matches no object. \\
FN & A ground-truth object remains unmatched after every prediction is evaluated. \\
\bottomrule
\end{tabularx}
\end{center}
\end{prioritybox}

\subsection{Type 5: Multiclass NMS}

After low-confidence predictions are removed, divide the remaining predictions by class and run NMS independently within each group.

\begin{formulabox}
\[
\boxed{\text{Run NMS separately for each class.}}
\]
\[
\boxed{\text{Boxes of different classes do not suppress one another.}}
\]
\end{formulabox}

Thus, a car prediction should not suppress a pedestrian prediction, even when their boxes overlap heavily.

\begin{practicebox}{Problem 1 --- Same location, different classes}
Consider
\[
A=(10,10,50,50),\quad \text{Car},\quad s_A=0.92,
\]
\[
B=(12,12,48,48),\quad \text{Car},\quad s_B=0.80,
\]
\[
C=(11,11,49,49),\quad \text{Pedestrian},\quad s_C=0.85,
\]
with \(T_{\mathrm{NMS}}=0.5\).

\textbf{Car class.} Compare \(A\) and \(B\). Their areas are
\[
A_A=40\times40=1600,
\qquad
A_B=36\times36=1296.
\]
Because \(B\) lies completely inside \(A\),
\[
\operatorname{IoU}(A,B)=\frac{1296}{1600}=0.81.
\]
Since \(0.81\geq0.5\), keep the higher-confidence car box \(A\) and suppress \(B\):
\[
\boxed{\text{Car survivor}=A}.
\]

\textbf{Pedestrian class.} Prediction \(C\) is the only pedestrian box, so
\[
\boxed{\text{Pedestrian survivor}=C}.
\]
Although \(C\) overlaps \(A\) strongly, the class labels differ. Therefore,
\[
\boxed{\text{Final detections}=\{A,C\}}.
\]
\end{practicebox}

\begin{practicebox}{Problem 2 --- Two classes with several boxes}
The prediction set is

\begin{center}
\begin{tabular}{ccc}
\toprule
\textbf{Box} & \textbf{Class} & \textbf{Confidence} \\
\midrule
\(A\) & Car & \(0.95\) \\
\(B\) & Car & \(0.83\) \\
\(C\) & Car & \(0.70\) \\
\(D\) & Pedestrian & \(0.90\) \\
\(E\) & Pedestrian & \(0.76\) \\
\bottomrule
\end{tabular}
\end{center}

Suppose
\[
\operatorname{IoU}(A,B)=0.68,
\qquad
\operatorname{IoU}(A,C)=0.20,
\]
\[
\operatorname{IoU}(D,E)=0.72,
\qquad
T_{\mathrm{NMS}}=0.5.
\]

\textbf{Car NMS.} Keep \(A\). Suppress \(B\) because \(0.68\geq0.5\), while \(C\) survives because \(0.20<0.5\). Thus,
\[
\boxed{\text{Car survivors}=\{A,C\}}.
\]

\textbf{Pedestrian NMS.} Keep \(D\). Since \(0.72\geq0.5\), suppress \(E\). Thus,
\[
\boxed{\text{Pedestrian survivor}=D}.
\]
Merging the class-wise results gives
\[
\boxed{\text{Final detections}=\{A,C,D\}}.
\]
\end{practicebox}

\begin{practicebox}{Problem 3 --- Confidence threshold + multiclass NMS}
Use the following predictions:

\begin{center}
\begin{tabular}{ccc}
\toprule
\textbf{Box} & \textbf{Class} & \textbf{Score} \\
\midrule
\(A\) & Car & \(0.91\) \\
\(B\) & Car & \(0.82\) \\
\(C\) & Car & \(0.55\) \\
\(D\) & Pedestrian & \(0.88\) \\
\(E\) & Pedestrian & \(0.61\) \\
\(F\) & Pedestrian & \(0.45\) \\
\bottomrule
\end{tabular}
\end{center}

Let
\[
T_{\mathrm{conf}}=0.60,
\qquad
T_{\mathrm{NMS}}=0.5,
\]
with
\[
\operatorname{IoU}(A,B)=0.70,
\qquad
\operatorname{IoU}(D,E)=0.30.
\]

\textbf{Step 1 --- Confidence filtering.} Remove \(C\) and \(F\) because their scores are below \(0.60\). The remaining set is
\[
\{A,B,D,E\}.
\]

\textbf{Step 2 --- Car NMS.} Keep \(A\), then suppress \(B\) because
\[
\operatorname{IoU}(A,B)=0.70\geq0.5.
\]
Thus, the car result is \(\{A\}\).

\textbf{Step 3 --- Pedestrian NMS.} Keep \(D\). Since
\[
\operatorname{IoU}(D,E)=0.30<0.5,
\]
\(E\) is not suppressed. Thus, the pedestrian result is \(\{D,E\}\).

The final result is
\[
\boxed{\text{Final detections}=\{A,D,E\}}.
\]
\end{practicebox}

\begin{practicebox}{Problem 4 --- Important cross-class trap}
Suppose
\[
A:\ \text{Car},\quad s_A=0.95,
\qquad
B:\ \text{Pedestrian},\quad s_B=0.90,
\]
and
\[
\operatorname{IoU}(A,B)=0.90,
\qquad
T_{\mathrm{NMS}}=0.5.
\]
Should \(B\) be removed? No:
\[
\boxed{B\text{ survives}}.
\]
The IoU value \(0.90\) would trigger suppression only if the two predictions were competing within the same class. Here, one is a car and the other is a pedestrian, so their cross-class IoU is irrelevant to class-wise NMS.
\end{practicebox}

\begin{practicebox}{Problem 5 --- Full coordinate example}
Consider
\[
A=(10,10,50,50),\quad \text{Car},\quad s_A=0.94,
\]
\[
B=(15,15,55,55),\quad \text{Car},\quad s_B=0.82,
\]
\[
C=(12,12,48,48),\quad \text{Pedestrian},\quad s_C=0.88,
\]
\[
D=(80,80,120,120),\quad \text{Pedestrian},\quad s_D=0.75,
\]
with \(T_{\mathrm{NMS}}=0.5\).

\textbf{Car class.} Boxes \(A\) and \(B\) each have area \(40^2=1600\). Their intersection is \(35\times35=1225\), so
\[
\operatorname{IoU}(A,B)
=\frac{1225}{1600+1600-1225}
=\frac{1225}{1975}
\approx0.620.
\]
Keep \(A\) and suppress \(B\).

\textbf{Pedestrian class.} Boxes \(C\) and \(D\) do not overlap, so both survive. Although \(C\) overlaps the car box \(A\), that cross-class overlap does not cause suppression.

Therefore,
\[
\boxed{\text{Final detections}=\{A,C,D\}}.
\]
\end{practicebox}

\subsubsection{Multiclass NMS exam recipe}

\begin{focusbox}[Use this order every time]
\begin{center}
\large\bfseries\color{PlanNavy}
Confidence Filtering
\(\boldsymbol{\longrightarrow}\)
Split by Class
\(\boldsymbol{\longrightarrow}\)
NMS Within Each Class
\(\boldsymbol{\longrightarrow}\)
Merge Survivors
\end{center}
\end{focusbox}

\begin{enumerate}[label=\textbf{\arabic*.},leftmargin=*,align=left,nosep]
  \item Remove every prediction with \(s<T_{\mathrm{conf}}\).
  \item Group the remaining boxes by predicted class.
  \item Sort each class group from highest to lowest confidence.
  \item Run the complete iterative NMS procedure inside each group.
  \item Combine the surviving boxes from all classes.
\end{enumerate}

\begin{prioritybox}{Mistake to avoid}
\[
\boxed{
\text{Never suppress a box merely because a different class overlaps it.}
}
\]
IoU-based suppression compares competing predictions for the same class; it does not remove valid detections of different object categories at the same location.
\end{prioritybox}

\subsection{Type 6: Anchor boxes and YOLO dimensions}

Practice:

\begin{itemize}
  \item assigning the anchor with the highest IoU;
  \item handling two objects whose centers fall in one grid cell;
  \item recognizing poor anchor matching;
  \item calculating the complete output tensor.
\end{itemize}

\begin{formulabox}
\[
\boxed{\text{YOLO output shape}=S\times S\times B(5+C)}
\]
Here \(S\times S\) is the grid, \(B\) is the number of predicted boxes or anchors per cell, and \(C\) is the number of classes.
\end{formulabox}

\newpage
\subsection{Type 7: Architecture answers}

Keep these answers short and exam-ready.

\begin{center}
\small
\begin{tabularx}{\textwidth}{>{\bfseries}p{2.9cm} X}
\toprule
Model & Core idea \\
\midrule
R-CNN & Generate proposals, then run the CNN separately on every proposal. \\
Fast R-CNN & Compute one shared CNN feature map, but retain external region proposals. \\
Faster R-CNN & Use a shared feature map and replace external proposals with a learned region proposal network (RPN). \\
YOLO & Perform one-stage, grid-based object detection in a single network pass. \\
U-Net & Use an encoder--decoder segmentation architecture with skip connections that recover spatial detail. \\
\bottomrule
\end{tabularx}
\end{center}

\subsection{Four hard detection challenges}

\begin{enumerate}[label=\textbf{Challenge \arabic*:},leftmargin=*,align=left,nosep]
  \item Six boxes, two classes, confidence filtering, and class-wise NMS.
  \item Five boxes and two ground-truth objects: perform NMS, then classify the survivors as TP or FP.
  \item Repeat one NMS problem with different thresholds and explain why the survivor set changes.
  \item Combine a YOLO grid, anchors, output dimensions, and anchor assignment.
\end{enumerate}

% -----------------------------------------------------------------------------
\section{Final Phase: Three Mock Finals}

Once the topic drills are complete, stop practicing topics in isolation and switch to timed, mixed papers.

\begin{center}
\small
\begin{tabularx}{\textwidth}{>{\bfseries}p{2.4cm} p{2.8cm} X}
\toprule
Paper & Level & Required mix \\
\midrule
Mock 1 & Standard & ANN forward/backprop; a CNN parameter count; IoU/NMS; two or three architecture concepts. \\
Mock 2 & Variation-heavy & Tanh + Softmax ANN; depthwise-separable CNN; confidence threshold + NMS + TP/FP; one architecture numerical. \\
Mock 3 & Hardest & Mixed-activation deep ANN; a full CNN with flattening and a depthwise block; multiclass NMS and anchors; short conceptual comparisons. \\
\bottomrule
\end{tabularx}
\end{center}

\subsection{Execution checklist for every problem}

\begin{multicols}{2}
\begin{itemize}[label=\checkbox,nosep]
  \item Identify the problem family.
  \item Write the relevant formula first.
  \item Mark all dimensions or box coordinates.
  \item Calculate in a fixed, visible order.
  \item Check signs, transposes, and floor operations.
  \item State the final answer with units or shape.
  \item For NMS, state the class and threshold.
  \item For gradients, verify \(dW\) matches \(W\).
\end{itemize}
\end{multicols}

\begin{tcolorbox}[
  enhanced,
  colback=PlanGold!8,
  colframe=PlanGold!80!black,
  boxrule=0.8pt,
  arc=2mm,
  left=2mm,right=2mm,top=1.4mm,bottom=1.4mm
]
\centering
\textbf{\color{PlanNavy}Finish line: roughly 35 serious drill problems + 3 timed mock exams.}\\[-0.1em]
\(\boxed{\text{see the question}\ \longrightarrow\ \text{identify the pattern}\ \longrightarrow\ \text{execute it quickly}}\)
\end{tcolorbox}

\end{document}
